\chapter{Introduction}
\label{ch:intro}
% A useful shape: open with the puzzle and its answer in one page, then only the background a reader needs
% for your data chapters, each block ending in the gap that leads to the next. Close with one aim per chapter.

\section{Overview}
\label{sec:overview}
[In one page: the question, why it matters, and what this dissertation shows.]

\section{Background}
\label{sec:background}
[What is known — only what the reader needs for the chapters that follow.]

\section{The gap}
\label{sec:gap}
[What is not known, and why it matters.]

\section{Hypothesis and specific aims}
\label{sec:aims}
\begin{itemize}
  \item Aim 1 (Chapter~\ref{ch:study1}): [aim]
  \item Aim 2 (Chapter~\ref{ch:study2}): [aim]
  \item Aim 3 (Chapter~\ref{ch:study3}): [aim]
\end{itemize}

\section{Strategy and limitations}
\label{sec:strategy}
[The approach across chapters, and the caveats a committee will ask about, named up front.]
